\documentclass[11pt]{article}
\usepackage[margin=1in]{geometry}
\usepackage{amsmath,amssymb}
\usepackage{parskip}
\usepackage[colorlinks=true,linkcolor=black,urlcolor=blue!60!black]{hyperref}

\newcommand{\HH}{\mathbb{H}}
\newcommand{\CC}{\mathbb{C}}
\newcommand{\tk}{\widetilde{k}_T}
\newcommand{\dT}{\delta_T}

\title{\vspace{-2.5em}Basepoint independence without cohomology,\\[0.2em]
\large and why $c_f(0)=0$ is a subtraction, not a restriction}
\author{}
\date{}

\begin{document}
\maketitle
\vspace{-3.5em}

\section*{The one identity that does the work}

Define the linear functional on a two-segment finite contour,
\[
  L^*(f,s;\tau_0)=\int_{C_S(\tau_0)} f(\tau)\,\tau^{s-1}\,d\tau
  \;+\;\int_{C_T(\tau_0)} f(\tau)\,\tk(\tau,s)\,d\tau ,
\]
with the Hurwitz $T$-kernel built from two terms,
$\tk(\tau,s)=\zeta(1-s,\tau+1)-\varpi_k\,\zeta\!\big(s-(k-1),\tau+1\big)$,
where $\varpi_k$ is the constant $S$-fold weight phase ($\varpi_k=i^k$ in the
convention of arXiv:1806.09874; \emph{not} an $s$-dependent factor). The kernel
exists to satisfy a single difference equation under
$\dT g(\tau):=g(\tau+1)-g(\tau)$:
\[
  \dT\,\tk(\tau,s)\;=\;-(\tau+1)^{s-1}+\varpi_k\,(\tau+1)^{\,k-1-s}.
\]
The two terms of $\tk$ are present precisely to produce \emph{both} powers,
$\tau^{s-1}$ and $\tau^{k-1-s}$ --- enough to cancel both the direct and the
$S$-image boundary terms below. (Exact prefactors are contour-specific; read them
off the actual geometry.)

\section*{$\dT$-inversion \emph{is} basepoint independence}

These are not two facts. Vary the base point and differentiate. The $T$-segment
endpoints move together, and the only property of $f$ used is $T$-periodicity
$f(\tau_0+1)=f(\tau_0)$, so the $T$-boundary variation is
\[
  \partial_{\tau_0}\!\!\int_{C_T}\! f\,\tk
  \;=\; f(\tau_0)\,\dT\tk(\tau_0,s).
\]
This is engineered to cancel the $S$-segment variation mode by mode. Every Fourier
mode of $f$ cancels except the constant, leaving
\[
  \boxed{\;\partial_{\tau_0} L^*(f,s;\tau_0)\;=\;c_f(0)\,\Phi(\tau_0,s)\;}
\]
with $\Phi$ \emph{universal} --- independent of $f$, of its principal part, of
everything but the single coefficient $c_f(0)$. So ``$\tk$ inverts $\dT$'' and
``$L^*$ is $\tau_0$-independent'' are the same statement. The cohomology was only
the name; the difference equation is the content, and it survives stripping the
cohomology away.

\section*{$c_f(0)=0$ is a subtraction, not a restriction}

The boxed identity says the entire base-point dependence is the one constant mode.
Hence, by linearity,
\[
  L^*\!\big(f-c_f(0),\,s;\tau_0\big)=L^*(f,s;\tau_0)-c_f(0)\,L^*(1,s;\tau_0)
\]
has $\partial_{\tau_0}=c_f(0)\Phi-c_f(0)\Phi=0$. Therefore \textbf{every} meromorphic
modular form --- any pole structure, constant term present or not --- has a
base-point-independent $L^*$ once the constant mode is removed. The subtracted piece
$c_f(0)\,L^*(1,s;\tau_0)$ is explicit, and is exactly where the poles of $\Lambda$
sit. Nothing is discarded; the constant is relocated to the polar part.

\paragraph{Precision (do not skip when writing the proof).}
$f-c_f(0)$ is \emph{not} a modular form: subtracting a weight-$0$ constant from a
weight-$k$ form breaks the $S$-law. So the independence is \emph{not} justified by
calling $f-c_f(0)$ ``the cusp form.'' It is justified by linearity together with the
universal-mode identity above --- i.e.\ by $\dT$-inversion. What the cancellation
actually uses, and what \emph{is} preserved under the subtraction, is
$T$-periodicity. Keep those two straight: $T$-periodicity survives the subtraction
and is used; $S$-modularity is broken by it and is not what carries the argument.

\section*{The poles come from the kernel, not from a divergence}

Classically, with $\Lambda(f,s)=\int_0^\infty\!\big(f(iy)-c_0\big)y^{s-1}dy$, one splits
at $y=1$ and folds the lower piece by $f(i/y)=(iy)^k f(iy)$; the $(f-c_0)$ integrals
are entire and the leftover elementary integrals of the constant give
\[
  \Lambda(f,s)=(\text{entire})-\frac{c_0}{s}-\frac{i^k c_0}{k-s},
\]
a simple pole at $s=0$ with residue $-c_0$ and one at $s=k$ with residue $i^k c_0$.
These sit on top of the $n\neq0$ incomplete-gamma sum without touching it: the BFK
part is entire, the poles are purely the $n=0$ mode.

In the finite contour there is no cusp, so the poles cannot come from a divergence at
$y\to0,\infty$. They come from the kernel. The Hurwitz function $\zeta(\sigma,a)$ has
its only pole at $\sigma=1$, with residue $1$ independent of $a$. The first kernel
term $\zeta(1-s,\tau+1)$ meets it at $s=0$; the second Hurwitz term, carrying the
$S$-fold weight phase, meets it at $s=k$. Near $s=0$,
\[
  \int_{C_T} f(\tau)\,\tk(\tau,s)\,d\tau
  \;\sim\; -\frac1s\int_{\tau_0}^{\tau_0+1} f(\tau)\,d\tau ,
\]
and the unit-length $T$-segment annihilates every nonzero Fourier mode,
$\int_{\tau_0}^{\tau_0+1} e^{2\pi i n\tau}\,d\tau=0$ for $n\neq0$, leaving exactly
$c_0$. So the residue at $s=0$ is $-c_0$, and it is \emph{basepoint-independent} even
though the regular part of $c_0\,L^*(1;\tau_0)$ is not. The $s=k$ pole arises the same
way through the second term; its residue is $i^k c_0$.

This matches McGady, arXiv:1806.09874, Thm.~1.2:
\[
  L_f^*(s)^{\mathrm{reg}}
  =-c_f(0)\!\left(\frac{B^s}{s}+\frac{i^k B^{k-s}}{k-s}\right)+\mathrm{regular}(s)
  = i^k L_f^*(k-s)^{\mathrm{reg}},
\]
so residue $-c_f(0)$ at $s=0$, residue $+i^k c_f(0)$ at $s=k$, and functional equation
$L^*(s)=i^k L^*(k-s)$, for every weight-$k$ meromorphic form regular at cusps.

\paragraph{Origin of the $i^k$.} The phase is the \emph{$S$-fold} factor, weight-indexed,
not $s$-indexed. On the imaginary axis $\tau=iy$, modularity $f(-1/\tau)=\tau^k f(\tau)$
folds the contour with $S(i/y)=iy$ and pulls out $(1/i)^k=(-i)^k$ --- a constant in $s$
depending only on the weight. That is why it survives as $i^k$ for \emph{all} even $k$
and attaches to the $s=k$ pole rather than the $s=0$ one. Any $s$-dependent phase of
the form $e^{i\pi(s-1)}$ is wrong here: at integer $s$ it is $\pm1$ and can never
produce a factor of $i$. The second kernel term's prefactor must be the constant
$S$-fold weight phase; take its exact form from the fold, not from an $s$-dependent
reconstruction.

\paragraph{Summary.}
The variation-in-$T$ operator is load-bearing, not an artifact. It is one operator
seen on two slices: on $V_{k-2}$ for the period polynomial (integer $s$), and on
functions with Hurwitz $\zeta$ as the antidifference of $\tau^{s-1}$ for the
$L$-function (all $s$). Its single difference equation is simultaneously the
parabolic gauge, the base-point independence, and the reason a constant-term
subtraction --- rather than a hypothesis on $f$ --- suffices for arbitrary pole
structure.

\end{document}
